% ==========================================================
% 020_motivation.tex
% Forecast Admissibility Surface (FAS)
% ==========================================================

\section{Motivation}
\label{sec:fas_motivation}

Modern forecasting systems often respond to poor predictive behavior by
escalating modeling complexity: richer feature sets, non-linear learners,
hierarchical pooling, or ensemble strategies. While such techniques can be
effective in many regimes, they implicitly assume that the underlying demand
slice is structurally amenable to learning and control. In operational settings,
this assumption does not universally hold.

Certain slices of the forecast domain exhibit properties that are fundamentally
hostile to stable modeling. These include extreme intermittency with long runs of
zeros, highly irregular support with insufficient observational mass, and error
distributions dominated by infrequent but severe spikes. In such regimes,
baseline forecasts already reveal instability: absolute errors are heavy-tailed,
spike frequencies are elevated, and empirical performance varies erratically
across adjacent time windows. Importantly, these behaviors persist across model
classes, indicating that the issue is not one of model inadequacy but of slice
structure.

Without an explicit admissibility construct, forecasting pipelines treat all
slices as equally eligible for learning. This leads to several undesirable
outcomes. Models are trained on slices where signal is weak or nonexistent,
inflating variance and computational cost. Evaluation metrics oscillate as a
function of rare events rather than systematic behavior. Governance and
readiness-control layers receive noisy or contradictory signals, making policy
decisions appear unstable or arbitrary. In the worst cases, forecasts produced
for such slices convey a false sense of precision while masking structural
uncertainty.

A common failure mode in these settings is the belief that additional data or
longer training windows will eventually resolve instability. While increased
support can be beneficial, many operational slices remain hostile even at scale.
Demand may be driven by rare, exogenous events; executional constraints may
truncate observable behavior; or production realities may impose sharp
discontinuities that defeat smooth learning. In these cases, the act of
forecasting itself becomes the source of error propagation rather than its
remedy.

The Forecast Admissibility Surface (FAS) is motivated by the need to distinguish
between slices where forecasting effort is justified and slices where it is not.
By examining baseline-derived error anatomy prior to advanced modeling, FAS
provides a principled mechanism for identifying structural hostility early in
the pipeline. This enables forecasting systems to allocate modeling capacity
where it is likely to yield stable gains, while constraining or blocking slices
where learning and control are inherently unreliable.

In this sense, FAS formalizes a governance question that is often left implicit:
\emph{Should this slice be forecasted at all?} By making this decision explicit,
deterministic, and auditable, FAS reduces downstream noise, clarifies the
interpretation of evaluation metrics, and strengthens the overall integrity of
the Forecast Readiness Framework.
